\section{Construcción del Conjunto de Datos}

El conjunto de datos se construye a partir del material de video registrado en el
propio recinto, en lugar de reutilizar un conjunto público. Esta decisión responde
a que las condiciones que determinan el desempeño del sistema ---perspectiva de
montaje, rango dinámico de la cámara, formato y tipografía de las matrículas
locales, patrones de iluminación del sitio--- no están representadas en los
conjuntos de uso general. El costo de esta decisión es la necesidad de construir
el etiquetado desde cero, lo que motiva el diseño de un proceso de extracción y
generación asistida que se describe en este apartado.

El proceso no consiste en muestrear fotogramas al azar, sino en localizar
primero los intervalos en los que ocurre algo relevante, reconstruir después la
trayectoria completa de cada vehículo y conservar únicamente los fotogramas con
potencial de contener una placa legible. La Figura \ref{fig:pipeline_dataset}
presenta el proceso, organizado en dos pasadas, y la Tabla \ref{tab:parametros}
enumera los parámetros que lo gobiernan.

\begin{figure}[htbp]
    \centering
    \begin{tikzpicture}
        % *** Entrada ***
        \spData{entrada}{-2.6,1.15}{video archivado}

        % *** Etapas ***
        \spBlock{etapa1}{-2.6,0}{7.4cm}{purpura}
            {Muestreo disperso y puerta de movimiento}{Primera pasada}
        \spBlock{etapa2}{-2.6,-1.35}{7.4cm}{purpura}
            {Re-decodificación densa}{Segunda pasada}
        \spBlock{etapa3}{-2.6,-2.70}{7.4cm}{purpura}
            {Detección y seguimiento}{YOLO26 + ByteTrack}
        \spBlock{etapa4}{-2.6,-4.05}{7.4cm}{purpura}
            {Clasificación de trayectorias}{Zonas y línea virtual}
        \spBlock{etapa5}{-2.6,-5.40}{7.4cm}{purpura}
            {Selección de fotogramas}{Mejores candidatos por trayectoria}
        \spBlock{etapa6}{-2.6,-6.75}{7.4cm}{purpura}
            {Generación sintética y revisión}{Ampliación de dominio y etiquetado}

        % *** Artefactos ***
        \spData{d1}{4.5,0}{ventanas de actividad}
        \spData{d2}{4.5,-1.35}{fotogramas continuos}
        \spData{d3}{4.5,-2.70}{trayectorias con identidad}
        \spData{d4}{4.5,-4.05}{ingreso, salida o tránsito}
        \spData{d5}{4.5,-5.40}{candidatos de placa}
        \spData{d6}{4.5,-6.75}{conjunto anotado}

        % *** Flujo ***
        \draw[spflow] (entrada.south) -- (etapa1.north);
        \draw[spflow] (etapa1.south) -- (etapa2.north);
        \draw[spflow] (etapa2.south) -- (etapa3.north);
        \draw[spflow] (etapa3.south) -- (etapa4.north);
        \draw[spflow] (etapa4.south) -- (etapa5.north);
        \draw[spflow] (etapa5.south) -- (etapa6.north);

        \draw[spflow] (etapa1.east) -- (d1.west);
        \draw[spflow] (etapa2.east) -- (d2.west);
        \draw[spflow] (etapa3.east) -- (d3.west);
        \draw[spflow] (etapa4.east) -- (d4.west);
        \draw[spflow] (etapa5.east) -- (d5.west);
        \draw[spflow] (etapa6.east) -- (d6.west);

    \end{tikzpicture}
    \caption{Pipeline de construcción del conjunto de datos.}
    \label{fig:pipeline_dataset}
\end{figure}

\begin{longtable}{
        P{0.30}
        P{0.12}
        P{0.50}
    }
    \caption{Parámetros de configuración del extractor}
    \label{tab:parametros}\\

    \toprule
    \rowcolor{gray!20}
    \textbf{Parámetro} &
    \textbf{Valor} &
    \textbf{Función} \\
    \midrule
    \endfirsthead

    \multicolumn{3}{c}{\textit{Continuación de la Tabla
    \ref{tab:parametros}}} \\
    \toprule
    \rowcolor{gray!20}
    \textbf{Parámetro} &
    \textbf{Valor} &
    \textbf{Función} \\
    \midrule
    \endhead

    \texttt{scan\_stride} & 12 & Separación entre muestras en la primera pasada. \\
    \texttt{motion\_downscale} & 0.25 & Reducción aplicada antes de estimar el movimiento. \\
    \texttt{motion\_pixel\_threshold} & 25 & Diferencia de intensidad que cuenta como cambio. \\
    \texttt{motion\_threshold} & 0.008 & Fracción de píxeles cambiados para declarar actividad. \\
    \texttt{segment\_pad\_frames} & 45 & Margen añadido a cada ventana detectada. \\
    \texttt{segment\_merge\_gap} & 60 & Separación máxima para fusionar ventanas contiguas. \\
    \texttt{dense\_stride} & 1 & Procesamiento fotograma a fotograma en la segunda pasada. \\
    \texttt{grace\_frames} & 30 & Fotogramas de tolerancia antes de cerrar una trayectoria. \\
    \texttt{min\_track\_hits} & 3 & Observaciones mínimas para aceptar una trayectoria. \\
    \texttt{top\_k\_frames} & 10 & Fotogramas conservados por trayectoria. \\
    \texttt{vehicle\_conf} & 0.35 & Confianza mínima de detección del vehículo. \\
    \texttt{min\_area\_pixels} & 3000 & Área mínima de la caja del vehículo. \\
    \texttt{ocr\_min\_conf} & 0.60 & Confianza mínima para aceptar la lectura preliminar de la placa. \\
    \texttt{min\_plate\_px} & 40 & Altura mínima del recorte de placa. \\
    \texttt{min\_plate\_sharpness} & 15.0 & Nitidez mínima exigida al recorte de placa. \\
    \texttt{plate\_y\_min} / \texttt{plate\_y\_max} & 0.55 / 1.00 & Franja vertical esperada de la placa dentro del vehículo. \\*
    \multicolumn{3}{r}{\footnotesize\textit{Fuente: Elaboración propia}} \\*
    \bottomrule
\end{longtable}

\subsection{Calibración de zonas de interés y línea virtual}

La pertenencia de un vehículo al recinto no puede resolverse mediante una única
región rectangular: el tránsito que circula por la vía pública y los vehículos
que ingresan ocupan la misma banda de la imagen, de modo que la posición por sí
sola no discrimina. Por esa razón la geometría se define mediante polígonos en
coordenadas normalizadas ---independientes de la resolución de captura--- que
distinguen tres elementos:

\begin{itemize}
    \item \textbf{Zona de vía:} polígono que cubre la calzada pública.
    \item \textbf{Zona interior:} polígono estrictamente posterior al acceso, que
        representa el espacio dentro del recinto.
    \item \textbf{Línea virtual:} segmento orientado que separa ambas zonas y
        cuya intersección detecta el cruce.
\end{itemize}

El punto de referencia empleado para evaluar la pertenencia y el cruce es el
\textit{punto de contacto con el suelo} ---el centro del borde inferior de la caja
del vehículo--- y no su centroide. La elección es física: un vehículo toca la
calzada por la parte inferior de su silueta, por lo que ese punto es el único que
atraviesa la línea en el mismo instante en que el vehículo lo hace.

La calibración no se realiza por inspección visual del fotograma, sino sobre
evidencia. Se dispone de una herramienta de visualización de trayectorias que
ejecuta el seguimiento sobre un fragmento de video sin aplicar ningún
filtro de zona y superpone todas las trayectorias reconstruidas sobre un
fotograma de referencia. Los vehículos que ingresan y los que sólo transitan
producen dos familias de curvas visualmente distinguibles, de manera que los
polígonos y la línea se trazan a partir de la separación observada entre ambas
familias y no de una estimación previa.

Un aspecto particular del entorno de captura es la presencia de información
incrustada por el grabador: el equipo superpone una etiqueta de canal y una marca
temporal directamente sobre el video. Dado que el motor de reconocimiento de
texto interpreta esa información como si fuese una matrícula ---llegando a leer
el identificador del canal---, la configuración declara regiones normalizadas que
se enmascaran antes de invocar al reconocimiento. Esta corrección es condición
necesaria para la validez de las muestras extraídas.

\subsection{Detección de movimiento y segmentación de ventanas}

Procesar la totalidad de los fotogramas de un archivo de video de varias horas
resulta prohibitive, pero alimentar un seguidor con fotogramas dispersos tampoco
es viable: los algoritmos de seguimiento suponen continuidad entre fotogramas
consecutivos y pierden las trayectorias cuando esa continuidad no existe. La
solución adoptada consiste en separar el problema en dos pasadas.

La primera pasada recorre el video de forma dispersa y calcula, en cada muestra,
una señal de movimiento de costo casi nulo: el fotograma se reduce a un cuarto de
su tamaño, se convierte a escala de grises, se suaviza y se compara con la muestra
anterior. La fracción de píxeles cuya diferencia supera un umbral de intensidad
constituye la energía de movimiento. Cuando esa energía supera un valor mínimo, la
muestra se marca como actividad.

Los instantes marcados se agrupan en ventanas temporales, a las que se añade un
margen por ambos extremos ---para no truncar la aproximación del vehículo--- y que
se fusionan cuando su separación es menor que un umbral, evitando fragmentar un
mismo tránsito en varios segmentos.

La segunda pasada vuelve a decodificar únicamente las ventanas
identificadas, esta vez fotograma a fotograma. Sólo en esta pasada se ejecutan la
detección y el seguimiento, sobre una secuencia continua. El resultado es una
reducción drástica del cómputo: el trabajo pesado se concentra en la fracción del
video donde efectivamente hay vehículos.

\subsection{Selección de fotogramas representativos}

El extractor razona sobre la trayectoria completa de cada vehículo y nunca sobre
un fotograma aislado. Esta decisión es la que permite distinguir un ingreso de un
simple tránsito sin depender de una única imagen: un vehículo que ingresa cruza la
línea virtual en dirección al recinto y su desplazamiento es predominantemente en
profundidad, con un aumento progresivo de su escala; un vehículo que sólo transita
mantiene una dirección horizontal y nunca alcanza la zona interior.

Las trayectorias con menos observaciones que el mínimo establecido se descartan
por considerarse ruido de detección. Para las restantes se conservan únicamente
los mejores fotogramas, seleccionados mediante un indicador económico de
legibilidad de placa. Ese indicador no evalúa el vehículo completo, sino su parte
inferior ---donde se ubica la placa---, donde mide la nitidez mediante la varianza
del laplaciano, y pondera el resultado por el tamaño del vehículo. Los recortes
que ocupan casi la totalidad del encuadre se penalizan deliberadamente, porque
corresponden a vehículos tan próximos que su placa queda fuera de cuadro o
recortada.

Sobre los fotogramas así elegidos se aplican además restricciones de calidad: se
exige una altura mínima del recorte de placa, una nitidez mínima y que la placa se
encuentre dentro de la franja vertical esperada. Los fotogramas que no satisfacen
estas condiciones no se descartan de inmediato, sino que se derivan al proceso de
revisión descrito a continuación.

\subsection{Pseudo-etiquetado, generación sintética y etiquetado humano}

Para garantizar la fiabilidad del entrenamiento y evitar la circularidad
metodológica, el etiquetado se organiza bajo reglas estrictas según el tipo de
atributo:

\begin{enumerate}
    \item \textbf{Pre-etiquetado y revisión de matrículas:} El motor PaddleOCR
        actúa exclusivamente como pseudo-etiquetador preliminar. Sus detecciones
        pasan a una carpeta de revisión (\texttt{review/}) para ser verificadas o
        corregidas por un humano antes de ingresar al conjunto de datos oficial.
        Para filtrar ruido antes de la revisión, se descartan cadenas sin dígitos
        (que corresponden a nombres de marcas impresas en la carrocería) y recortes
        con relaciones de aspecto ajenas al estándar rectangular.

    \item \textbf{Etiquetado humano directo de atributos (tipo y color):} Los
        atributos visuales del vehículo (categoría de carrocería y color
        predominante) son etiquetados en su totalidad por un operador humano mediante
        una herramienta interactiva especializada (\texttt{label\_type.py}).
        Se prohíbe explícitamente el uso de clasificadores automáticos para generar
        estas etiquetas: hacerlo implicaría entrenar el modelo final con las
        predicciones de una regla contra la cual se pretende evaluar su desempeño,
        incurriendo en una invalidez científica.

    \item \textbf{Generación sintética de matrículas (pre-entrenamiento):} Debido
        a que la recolección de placas reales requiere procesar cientos de horas de
        video, se implementa un generador sintético
        (\texttt{generate\_synth\_plates.py}). Este módulo renderiza la tipografía,
        colores y símbolos oficiales de las placas vehiculares de Bolivia (formatos
        \texttt{NNNNLLL} y \texttt{NNNLLL} con leyenda nacional) y aplica un
        \textit{pipeline} de degradación física (deformación de perspectiva,
        variación fotométrica, sombras sintéticas, desenfoque de movimiento y
        ruido de sensor) ajustando su tamaño a la distribución real observada
        (mediana de $96 \times 39$ píxeles). Las muestras sintéticas se identifican
        en el manifiesto con la marca \texttt{synth: true} y se utilizan
        exclusivamente durante la fase de pre-entrenamiento, jamás para la
        evaluación final.
\end{enumerate}

La salida del proceso se organiza según la Tabla \ref{tab:salida_dataset}.

\begin{longtable}{
        P{0.30}
        P{0.65}
    }
    \caption{Organización de la salida del extractor y generador}
    \label{tab:salida_dataset}\\

    \toprule
    \rowcolor{gray!20}
    \textbf{Salida} &
    \textbf{Contenido} \\
    \midrule
    \endfirsthead

    \multicolumn{2}{c}{\textit{Continuación de la Tabla
    \ref{tab:salida_dataset}}} \\
    \toprule
    \rowcolor{gray!20}
    \textbf{Salida} &
    \textbf{Contenido} \\
    \midrule
    \endhead

    \texttt{vehicles/images/} & Recorte del vehículo con la placa visible. \\
    \midrule
    \texttt{vehicles/labels/} & Caja delimitadora de la placa en coordenadas del recorte. \\
    \midrule
    \texttt{plates/images/} & Recorte ajustado de la placa, destinado al OCR. \\
    \midrule
    \texttt{plates/labels/} & Texto de la matrícula confirmado por el operador. \\
    \midrule
    \texttt{attributes/labels.jsonl} & Etiquetas de tipo y color asignadas por el operador. \\
    \midrule
    \texttt{review/} & Lecturas de baja confianza o mala nitidez, pendientes de revisión. \\
    \midrule
    \texttt{negatives/} & Recortes de vehículos que transitan por la vía sin ingresar. \\
    \midrule
    \texttt{manifest.jsonl} & Registro central de muestras (etiqueta real o \texttt{synth: true}). \\*
    \multicolumn{2}{r}{\footnotesize\textit{Fuente: Elaboración propia}} \\*
    \bottomrule
\end{longtable}

\subsection{Participación, composición y protocolo de ablación del conjunto de datos}

La partición entre entrenamiento, validación y prueba se establece por sesión de
grabación, agrupada en bloques horarios, y no de forma aleatoria por imagen. Un
mismo vehículo aparece en decenas de fotogramas consecutivos; si esos fotogramas
se distribuyeran al azar, el modelo memorizaría el vehículo específico en lugar de
aprender a generalizar. Al mantener cada sesión en un único subconjunto, ninguna
trayectoria queda repartida entre entrenamiento y prueba, de modo que se elimina
la fuga de información por identidad.

Para evaluar el impacto de la generación sintética sobre el reconocimiento de
matrículas, se construye un protocolo de ablación determinista
(\texttt{prepare\_ocr\_ablation.py}) fijando una semilla aleatoria ($seed = 7$) y una
proporción de prueba del $20\%$.

La Tabla \ref{tab:caracterizacion_dataset} resume la composición del conjunto de
datos de reconocimiento de matrículas. La composición completa de los conjuntos de
atributos y de placas se detalla en el Anexo~B.

\begin{longtable}{
        P{0.20}
        P{0.12}
        P{0.60}
    }
    \caption{Composición del conjunto de datos de reconocimiento de matrículas}
    \label{tab:caracterizacion_dataset}\\

    \toprule
    \rowcolor{gray!20}
    \textbf{Conjunto} &
    \textbf{Muestras} &
    \textbf{Propósito y origen} \\
    \midrule
    \endfirsthead

    \multicolumn{3}{c}{\textit{Continuación de la Tabla
    \ref{tab:caracterizacion_dataset}}} \\
    \toprule
    \rowcolor{gray!20}
    \textbf{Conjunto} &
    \textbf{Muestras} &
    \textbf{Propósito y origen} \\
    \midrule
    \endhead

    \texttt{train\_real} & 1,085 & Muestras reales extraídas de las sesiones de grabación. \\
    \midrule
    \texttt{train\_synth} & 360,000 & Muestras sintéticas generadas con formato boliviano. \\
    \midrule
    \texttt{test} (evaluación) & 152 & Muestras reales independientes para la prueba final. \\
    \midrule
    \textbf{Total real procesado} & \textbf{1,237} & \textbf{$640$ placas distintas, particionadas por cadena de matrícula.} \\*
    \multicolumn{3}{r}{\footnotesize\textit{Fuente: Elaboración propia}} \\*
    \bottomrule
\end{longtable}

Además de la partición de entrenamiento y prueba del reconocedor, la evaluación de
extremo a extremo se realizó sobre $99$ recortes de vehículo con texto de placa
verificado por un operador, provenientes de los recortes del conjunto de prueba que
conservan la imagen original del vehículo. Ese conjunto no participa del
entrenamiento y se describe en el Anexo~B. El conjunto de color, por su parte,
excluye los minibuses y la clase residual \textit{otro}, de modo que su tamaño no
coincide con el del conjunto de tipo; la composición y el motivo de cada exclusión
se detallan en el mismo anexo.

Dos limitaciones del material condicionan el alcance de los resultados. La primera
es el volumen de matrículas reales: $640$ confirmadas por un operador, de las
cuales $541$ placas se destinan al ajuste fino, constituyen la restricción
dominante del módulo de reconocimiento y no la arquitectura del modelo
(Sección~\ref{sec:resultados_ocr}). Los protocolos de extracción están diseñados
para escalar la recolección sin alterar la metodología ni los criterios de
evaluación. La segunda es la procedencia: todas las sesiones provienen de las
mismas cámaras del recinto, de modo que la partición por sesión controla la fuga
por identidad dentro de una misma toma, y la partición por cadena de matrícula
elimina la fuga de placa, pero no se controla la variabilidad entre modelos de
cámara o condiciones de instalación distintas.

La integridad del material se verificó antes de la partición. Durante la
transferencia desde el disco externo de trabajo ---que presentó errores de
entrada/salida en aproximadamente el $1{.}3\%$ de los recortes del pool--- ocho
archivos sintéticos quedaron irrecuperables y otros veinticinco se truncaron. Los
treinta y tres se regeneraron de forma determinística: el generador deriva el
estado aleatorio de cada placa de su índice, de modo que la misma revisión del
código reproduce el mismo espécimen. En los veinticinco casos con etiqueta previa
se verificó que el texto regenerado coincidiera exactamente con el original; en
los ocho restantes la revisión se confirmó por continuidad con índices vecinos
recuperados. El conjunto final conserva $360{,}000$ especímenes válidos.

El episodio se documenta porque un generador determinístico por índice solo es
reproducible para una revisión fija del código. Una modificación posterior del
alfabeto de caracteres hace que un mismo índice produzca una placa distinta, y por
lo tanto que reanudar una generación interrumpida reescriba el material en vez de
continuarlo.
